% !TEX program = lualatex
% !BIB program = biber
%                                                __
%  __                                     __    /\ \
% /\_\    ___      __          __   _ __ /\_\   \_\ \
% \/\ \ /' _ `\  /'_ `\      /'_ `\/\`'__\/\ \  /'_` \
%  \ \ \/\ \/\ \/\ \L\ \  __/\ \L\ \ \ \/ \ \ \/\ \L\ \
%   \ \_\ \_\ \_\ \____ \/\_\ \____ \ \_\  \ \_\ \___,_\
%    \/_/\/_/\/_/\/___L\ \/_/\/___L\ \/_/   \/_/\/__,_ /
%                  /\____/     /\____/
%                  \_/__/      \_/__/
%
%                        journal ing.grid
%
%
% Template for authors
%
% Structure of this template:
% .
% ├── acknowledgements.tex <-- acknowledgements, funding information
% ├── figures              <-- images and pictures, same in all languages
% │   └── inggrid_Dashboard.png <- e.g. the dashboard on page 12
% ├── final.def            <-- custom title page and metadata fields
% ├── fonts                <-- needed fonts
% ├── comic_files          <-- everything that is the same in all languages
% │   ├── BLANKO_1-12_v1.0.0.pdf   <- blank comic pages, the drawings
% │   └── pages_functions.tex      <- comic macros, e.g. \comic{deu}
% ├── inggrid.cls          <-- the documentclass
% ├── language_files       <-- one folder per comic language
% │   ├── deu                  <- everything of the German version
% │   │   ├── deu.tex              <- all comic texts in German
% │   │   ├── FDI_deu.png          <- language dependent images
% │   │   └── pages_deu            <- text placement of the German pages
% │   │       ├── page_1.tex           <- one file per comic page
% │   │       └── ...
% │   ├── qya                  <- Quenya (Elvish) version, same structure, plus
% │   │   ├── fonts_qya.tex        <- language specific font declarations
% │   │   └── fonts_qya            <- language specific font files
% │   ├── eng                  <- English version, same structure
% │   └── lat                  <- Latin version, same structure
% ├── License              <-- the license file
% ├── logos                <-- needed logos
% ├── main.tex             <-- this file, content
% ├── metadata             <-- metadata of article
% │   └── article.tex          <- example article metadata (changed only by editors)
% │   └── authors.tex          <- example author data
% ├── README.md            <-- instructions
% ├── references.bib       <-- the bibliography
% ├── scikgtex.lua         <-- script for managing scikgtex metadata
% └── scikgtex.sty         <-- provides commands for research contribution annotations
%
%
\documentclass{inggrid}

%
%
% The documentclass already loads a number of packages, amongst others:
% - amsmath,amssym
% - xcolor, graphicx
% - babel, csquotes
% - biblatex
% - fontspec, microtype
% - booktabs, longtable
% - listings
%
% -----------------------------------------------
\usepackage{graphicx}
\usepackage{tikz}
\DeclareRobustCommand{\TikZ}{Ti\textit{k}Z}
\usetikzlibrary{decorations.text}
\usetikzlibrary{shapes.geometric}
\usepackage{eso-pic}
\usepackage{fontspec}
\usepackage{transparent}
\usepackage[nolist,nohyperlinks]{acronym}
\definecolor{FDIBrown}{RGB}{96,45,1}


% ------------------------------------------------------------
% FDI language switch
% ------------------------------------------------------------
\newcommand{\FDILanguage}{eng}
% Select the primary comic edition here. Register its code and display name in
% \FDILanguageList and \FDISetLanguageName in comic_files/pages_functions.tex.
% Other registered editions are added to the navigation and appendix
% automatically by \FDIJumpToAll and \FDIAppendixComics; they require no
% further change in this file.

% ------------------------------------------------------------
% FDI setup
% ------------------------------------------------------------

\pagestyle{empty}
\setlength{\parindent}{0pt}

\newcounter{fdipage}

% ------------------------------------------------------------
% Comic folder and blank comic pages
% ------------------------------------------------------------
% Enthaelt diese Datei und das Blanko-PDF, also alles, was fuer
% alle Sprachen gleich ist. Bei einer neuen Blanko-Version nur
% \FDIBlankoFile anpassen.

\newcommand{\FDIComicFolder}{comic_files}
\newcommand{\FDIBlankoFile}{\FDIComicFolder/BLANKO_1-12_v1.0.0.pdf}

% ------------------------------------------------------------
% Fonts setup depending on language
% ------------------------------------------------------------

% Normale Fonts
\newfontfamily{\FDIKarolingischeMinuskelDefault}{Karolingische_Minuskel.ttf}[
  Path = fonts/,
  FakeBold = 2
]

\newfontfamily{\FDIOverlayFontDefault}{KOMTXT__.ttf}[
  Path = fonts/
]

\newfontfamily{\FDIArialBOLDFontDefault}{ArialCEMTBlack.ttf}[
  Path = fonts/
]

\newfontfamily{\FDIArialFontDefault}{ARIAL.ttf}[
  Path = fonts/
]

% Elbische FDI-Schriften
% Wichtig: Diese Fonts werden NICHT global auf das Dokument angewendet,
% sondern nur über die FDI-Font-Makros unten verwendet.
% Renderer=Basic ist bei alten Tengwar-/Symbolfonts oft wichtig.

% Normale FDI-Sprechblasen / Standard-Overlay
\newfontfamily{\FDIElbArialFont}{tngani.ttf}[
  Path = fonts/,
  Renderer = Basic,
  Scale = 0.78
]

% Fette FDI-Sprechblasen / Hervorhebungen
\newfontfamily{\FDIElbArialBOLDFont}{tngani.ttf}[
  Path = fonts/,
  Renderer = Basic,
  Scale = 0.78,
  FakeBold = 1.8
]

% FDI-Overlay-Schrift, etwas größer für Comic-Texte
\newfontfamily{\FDIElbOverlayFont}{tngani.ttf}[
  Path = fonts/,
  Renderer = Basic,
  Scale = 0.82
]

% FDI-Karolingische-Minuskel-Ersatzschrift
\newfontfamily{\FDIElbKarolingischeMinuskel}{tngani.ttf}[
  Path = fonts/,
  Renderer = Basic,
  Scale = 0.86,
  FakeBold = 1.3
]

% ------------------------------------------------------------
% Public font aliases
% ------------------------------------------------------------

\providecommand{\KarolingischeMinuskel}{}
\providecommand{\FDIOverlayFont}{}
\providecommand{\ArialBOLDFont}{}
\providecommand{\ArialFont}{}

% ------------------------------------------------------------
% Apply FDI font aliases depending on language
% ------------------------------------------------------------
% Wichtig:
% Dieser Schalter verändert NICHT \setmainfont, \setsansfont oder \setmonofont.
% Er betrifft ausschließlich die FDI-Font-Makros:
% \KarolingischeMinuskel, \FDIOverlayFont, \ArialBOLDFont und \ArialFont.

\newcommand{\FDISetupLanguageFonts}{%
  \ifdefstring{\FDILanguage}{elb}{%
    \typeout{FDI: Language is elb. Switching FDI font macros only.}%

    % Nur FDI-Font-Makros auf elbische Varianten setzen.
    % Keine globalen Dokumentfonts ändern!
    \let\KarolingischeMinuskel\FDIElbKarolingischeMinuskel
    \let\FDIOverlayFont\FDIElbOverlayFont
    \let\ArialBOLDFont\FDIElbArialBOLDFont
    \let\ArialFont\FDIElbArialFont
  }{%
    \typeout{FDI: Language is not elb. Using default FDI font macros.}%

    % Normale FDI-Font-Makros setzen.
    \let\KarolingischeMinuskel\FDIKarolingischeMinuskelDefault
    \let\FDIOverlayFont\FDIOverlayFontDefault
    \let\ArialBOLDFont\FDIArialBOLDFontDefault
    \let\ArialFont\FDIArialFontDefault
  }%
}

% ------------------------------------------------------------
% Language folder depending on language
% Usage result: language_files/de, language_files/en, ...
% Enthaelt die Sprachdatei, den zugehoerigen pages-Ordner und
% alle sprachabhaengigen Bilder.
% ------------------------------------------------------------

\newcommand{\FDILanguageFolder}{%
  language_files/\FDILanguage%
}

% ------------------------------------------------------------
% Input language text file
% Loads: language_files/de/de.tex, language_files/en/en.tex, ...
% ------------------------------------------------------------

\newcommand{\FDIInputLanguageFile}{%
  \IfFileExists{\FDILanguageFolder/\FDILanguage.tex}{%
    \input{\FDILanguageFolder/\FDILanguage.tex}%
  }{%
    \PackageError{FDI}{Language file '\FDILanguageFolder/\FDILanguage.tex' not found}{%
      Check \string\FDILanguage\space and create the corresponding file.%
    }%
  }%
}

% ------------------------------------------------------------
% Set page folder depending on language
% Usage result: language_files/de/pages_de, language_files/en/pages_en, ...
% ------------------------------------------------------------

\newcommand{\FDIPageFolder}{%
  \FDILanguageFolder/pages_\FDILanguage%
}

% ------------------------------------------------------------
% Input page by number
% Usage: \FDIInputPage{1}
% Loads: language_files/de/pages_de/page_1.tex, ...
% ------------------------------------------------------------

\newcommand{\FDIInputPage}[1]{%
  \IfFileExists{\FDIPageFolder/page_#1.tex}{%
    \input{\FDIPageFolder/page_#1.tex}%
  }{%
    \PackageError{FDI}{Page file '\FDIPageFolder/page_#1.tex' not found}{%
      Check \string\FDILanguage\space and create the corresponding page file.%
    }%
  }%
}

% ------------------------------------------------------------
% Whole comic in one language
% Usage: \comic{de}        alle Comicseiten auf Deutsch
%        \comic[3]{en}     nur die ersten 3 Seiten, also Seite 1 bis 3
%        \comic[3-5]{en}   Seite 3 bis 5
%        \comic[7-7]{en}   nur Seite 7
% ------------------------------------------------------------
% Gibt alle Comicseiten der Sprache #2 direkt hintereinander aus.
% Sprache, Fonts und Texte gelten nur innerhalb von \comic. Nach
% \comic gelten wieder die Einstellungen des Dokuments, also
% \FDILanguage sowie \title/\subtitle/\shorttitle aus der in der
% Praeambel geladenen Sprachdatei.
% Dadurch koennen auch mehrere Sprachen nacheinander ausgegeben
% werden, z.B.:
%   \comic{de}
%   \comic{en}

% Anzahl der Comicseiten pro Sprache
\newcommand{\FDIComicPageCount}{12}

\makeatletter

\newcounter{FDIcomicpage}

% Optionales Argument von \comic auswerten.
% Aufruf immer mit angehaengtem --\FDI@stop, dadurch gilt:
%   "5"    ->  #1=5  #2=leer  ->  Seite 1 bis 5
%   "3-5"  ->  #1=3  #2=5     ->  Seite 3 bis 5
\def\FDI@stop{}

\def\FDI@parserange#1-#2-#3\FDI@stop{%
  \def\FDI@tmp{#2}%
  \ifx\FDI@tmp\@empty
    \def\FDI@firstpage{1}%
    \edef\FDI@lastpage{#1}%
  \else
    \edef\FDI@firstpage{#1}%
    \edef\FDI@lastpage{#2}%
  \fi
}

\newcommand{\comic}[2][\FDIComicPageCount]{%
  \begingroup
    % \edef, damit auch \comic{\FDILanguage} funktioniert
    \edef\FDILanguage{#2}%
    \FDISetupLanguageFonts
    \FDIInputLanguageFile
    % Seitenbereich bestimmen
    \edef\FDI@tmp{#1}%
    \expandafter\FDI@parserange\FDI@tmp--\FDI@stop
    % Beide Zaehler auf die Seite vor der ersten gewuenschten Seite setzen.
    % fdipage waehlt die Seite im Hintergrund-PDF und laeuft deshalb mit,
    % damit \comic[3-5] auch die Zeichnungen der Seiten 3 bis 5 zeigt.
    \setcounter{FDIcomicpage}{\FDI@firstpage}%
    \addtocounter{FDIcomicpage}{-1}%
    \setcounter{fdipage}{\value{FDIcomicpage}}%
    \@whilenum\value{FDIcomicpage}<\FDI@lastpage\do{%
      \stepcounter{FDIcomicpage}%
      \expandafter\FDIInputPage\expandafter{\number\value{FDIcomicpage}}%
    }%
    % letzte Comicseite ausgeben, solange die Sprache noch gesetzt ist
    \clearpage
  \endgroup
}

\makeatother

% ------------------------------------------------------------
% FDI header/footer offsets
% ------------------------------------------------------------
% Positive X-Werte verschieben nach rechts.
% Negative X-Werte verschieben nach links.
% Positive Y-Werte verschieben nach oben.
% Negative Y-Werte verschieben nach unten.

\newlength{\FDIHeaderXOffset}
\newlength{\FDIHeaderYOffset}
\newlength{\FDIFooterXOffset}
\newlength{\FDIFooterYOffset}

\setlength{\FDIHeaderXOffset}{0pt}
\setlength{\FDIHeaderYOffset}{20pt}
\setlength{\FDIFooterXOffset}{0pt}
\setlength{\FDIFooterYOffset}{-35pt}

% ------------------------------------------------------------
% FDI pagestyle: original inggrid pagestyle with offsets
% ------------------------------------------------------------

\makeatletter

\let\FDI@original@ps@inggrid\ps@inggrid

\def\ps@inggridfdi{%
  \FDI@original@ps@inggrid%

  \let\FDI@oddhead\@oddhead%
  \let\FDI@evenhead\@evenhead%
  \let\FDI@oddfoot\@oddfoot%
  \let\FDI@evenfoot\@evenfoot%

  \def\@oddhead{%
    \hspace*{\FDIHeaderXOffset}%
    \raisebox{\FDIHeaderYOffset}[0pt][0pt]{\FDI@oddhead}%
    \hspace*{-\FDIHeaderXOffset}%
  }%

  \def\@evenhead{%
    \hspace*{\FDIHeaderXOffset}%
    \raisebox{\FDIHeaderYOffset}[0pt][0pt]{\FDI@evenhead}%
    \hspace*{-\FDIHeaderXOffset}%
  }%

  \def\@oddfoot{%
    \hspace*{\FDIFooterXOffset}%
    \raisebox{\FDIFooterYOffset}[0pt][0pt]{\FDI@oddfoot}%
    \hspace*{-\FDIFooterXOffset}%
  }%

  \def\@evenfoot{%
    \hspace*{\FDIFooterXOffset}%
    \raisebox{\FDIFooterYOffset}[0pt][0pt]{\FDI@evenfoot}%
    \hspace*{-\FDIFooterXOffset}%
  }%
}

\makeatother


% ------------------------------------------------------------
% FDI pages
% ------------------------------------------------------------

\newcommand{\FDIPage}[1]{%
  \clearpage
  \stepcounter{fdipage}%
  \thispagestyle{inggridfdi}%
  \AddToShipoutPictureBG*{%
    \begin{tikzpicture}[remember picture, overlay]
      \node[anchor=center, inner sep=0pt]
        at (current page.center) {%
          \includegraphics[
            page=\value{fdipage},
            width=\paperwidth,
            height=\paperheight
          ]{\FDIBlankoFile}%
        };
    \end{tikzpicture}%
  }%
  #1%
  \null
}

% ------------------------------------------------------------
% FDI Texts
% ------------------------------------------------------------
% Globale Variable für den Zeilenabstand
\newlength{\FDITextLineSkip}
\setlength{\FDITextLineSkip}{12pt}


% \FDIStyledText{schriftgröße}{ignored}{text}
% Achtung: kein äußerer {...}-Block, damit \baselineskip bis zum Absatzende wirkt.
\newcommand{\FDIStyledText}[3]{%
  \FDIOverlayFont
  \fontsize{#1}{\FDITextLineSkip}\selectfont
  \setlength{\baselineskip}{\FDITextLineSkip}%
  #3%
}

% \FDIText{x}{y}{breite}{text}
\newcommand{\FDIText}[4]{%
  \node[
    anchor=north west,
    inner sep=0pt,
    outer sep=0pt,
    minimum width=#3
  ] at ([xshift=#1,yshift=-#2]current page.north west) {%
    \begin{minipage}[t]{#3}
      \centering
      \setlength{\lineskip}{0pt}%
      \setlength{\parskip}{0pt}%
      \FDIOverlayFont
      #4%
    \end{minipage}%
  };
}

% \FDITextLeft{x}{y}{breite}{text}
\newcommand{\FDITextLeft}[4]{%
  \node[
    anchor=north west,
    inner sep=0pt,
    outer sep=0pt,
    minimum width=#3
  ] at ([xshift=#1,yshift=-#2]current page.north west) {%
    \begin{minipage}[t]{#3}
      \raggedright
      \setlength{\lineskip}{0pt}%
      \setlength{\parskip}{0pt}%
      \FDIOverlayFont
      #4%
    \end{minipage}%
  };
}

\newcommand{\FDIStyledTextMinuskel}[3]{%
  {\KarolingischeMinuskel\fontsize{#1}{#2}\selectfont #3}%
}

% \FDITextLS{x}{y}{breite}{schriftgröße}{zeilenabstand}{text}
% \FDIText{x}{y}{breite}{schriftgröße}{zeilenabstand}{text}
\newcommand{\FDITextLS}[6]{%
  \node[
    anchor=north west,
    text width=#3,
    align=center,
    font=\FDIOverlayFont,
    inner sep=0pt,
    outer sep=0pt,
    minimum width=#3,
    font=\FDIOverlayFont\fontsize{#4}{#5}\selectfont
  ] at ([xshift=#1,yshift=-#2]current page.north west) {#6};
}



% \FDIOvalWrapText{x}{y}{breite}{text}
\newcommand{\FDIOvalText}[4]{%
  % unsichtbare / randlose Sprechblasenfläche
  \node[
    ellipse,
    draw=none,
    minimum width=#3,
    minimum height=1.9cm,
    inner sep=0pt,
    outer sep=0pt
  ] at ([xshift=#1,yshift=-#2]current page.north west) {};

  % Text mit angenäherter ovaler Kontur
  \node[
    anchor=center,
    inner sep=0pt,
    outer sep=0pt
  ] at ([xshift=#1,yshift=-#2]current page.north west) {%
    \begin{minipage}{#3}
      \centering
      \FDIOverlayFont
      \parshape=5
        0.55cm \dimexpr#3-1.10cm\relax
        0.25cm \dimexpr#3-0.50cm\relax
        0.00cm #3
        0.25cm \dimexpr#3-0.50cm\relax
        0.55cm \dimexpr#3-1.10cm\relax
      #4%
    \end{minipage}%
  };
}

% \FDIArcTextOutside{x}{y}{start angle}{end angle}{radius}{text}{style}
% Text ist von außen lesbar.
\newcommand{\FDIArcTextOutside}[7]{%
  \path[
    decorate,
    decoration={
      text along path,
      text={|#7| #6},
      text align=center
    }
  ]
  ([xshift=#1,yshift=-#2]current page.north west)
  arc[start angle=#3,end angle=#4,radius=#5];
}

% \FDIArcTextInside{x}{y}{start angle}{end angle}{radius}{text}{style}
% Text ist von innen lesbar.
\newcommand{\FDIArcTextInside}[7]{%
  \path[
    decorate,
    decoration={
      text along path,
      reverse path,
      text={|#7| #6},
      text align=center
    }
  ]
  ([xshift=#1,yshift=-#2]current page.north west)
  arc[start angle=#3,end angle=#4,radius=#5];
}

% \FDITextRotated{x}{y}{breite}{winkel}{text}
% winkel: positiver Wert = gegen den Uhrzeigersinn, negativer Wert = im Uhrzeigersinn
\newcommand{\FDITextRotated}[5]{%
  \node[
    anchor=north west,
    text width=#3,
    align=center,
    inner sep=0pt,
    outer sep=0pt,
    rotate=#4
  ] at ([xshift=#1,yshift=-#2]current page.north west) {#5};
}

% ------------------------------------------------------------
% FDI Text Posistion Makros
% ------------------------------------------------------------
% \def statt \newcommand, damit dieselbe Sprachdatei mehrfach bzw.
% mehrere Sprachdateien nacheinander geladen werden koennen (siehe \comic).
\newcommand{\FDISetText}[4]{%
  \expandafter\def\csname s#1p#2t#3\endcsname{#4}%
}

\newcommand{\FDIUseText}[3]{%
  \csname s#1p#2t#3\endcsname
}

% ------------------------------------------------------------
% FDI city dots
% ------------------------------------------------------------
% \FDICityDot{x}{y}{größe}{füllfarbe}
% Setzt einen Punkt mit weißer Outline.
% Alle Werte werden vom oberen linken Seitenrand gemessen.
% Beispiel: \FDICityDot{7.2cm}{11.4cm}{7pt}{red!85!black}

\newlength{\FDICityDotOutlineWidth}
\setlength{\FDICityDotOutlineWidth}{1.2pt}

\tikzset{
  FDICityDotStyle/.style={
    circle,
    draw=white,
    line width=\FDICityDotOutlineWidth,
    inner sep=0pt,
    outer sep=0pt
  }
}

\newcommand{\FDICityDot}[4]{%
  \node[
    FDICityDotStyle,
    minimum size=#3,
    fill=#4
  ] at ([xshift=#1,yshift=-#2]current page.north west) {};
}


% ------------------------------------------------------------
% FDI movable line numbers / panel numbers
% Uses original inggrid / lineno counter
% ------------------------------------------------------------

\makeatletter

% Zentraler Stil der Nummern aus lineno / inggrid
\newcommand{\FDILineNumberStyle}{%
  \linenumberfont
}

% \FDILineNumberBox{x}{y}
% x = Abstand vom linken Seitenrand
% y = Abstand vom oberen Seitenrand
% Nutzt den originalen linenumber-Zähler fortlaufend weiter.
\newcommand{\FDILineNumberBox}[2]{%
  \stepcounter{linenumber}%
  \node[
    anchor=north west,
    inner sep=0pt,
    outer sep=0pt
  ] at ([xshift=#1,yshift=-#2]current page.north west) {%
    {\FDILineNumberStyle\thelinenumber}%
  };%
}

\makeatother


% ------------------------------------------------------------
% Horizontal lines
% ------------------------------------------------------------

\newcommand{\FDIHLine}[3]{%
  \draw[lightgray, line width=0.4pt]
    ([xshift=#1,yshift=-#3]current page.north west)
    --
    ([xshift=#2,yshift=-#3]current page.north west);
}


% ------------------------------------------------------------
% Point to point lines
% ------------------------------------------------------------

% \FDILine{x1}{y1}{x2}{y2}
% zieht eine Linie von Punkt 1 zu Punkt 2
% alle Werte werden von oben links gemessen
\newcommand{\FDILine}[4]{%
  \draw[black, line width=0.4pt]
    ([xshift=#1,yshift=-#2]current page.north west)
    --
    ([xshift=#3,yshift=-#4]current page.north west);
}

% Sprachspezifische Schriften aus den Sprachordnern laden,
% z.B. language_files/qya/fonts_qya.tex. Muss in die Praeambel.
\FDIInputLanguageFonts

\FDISetupLanguageFonts
\FDIInputLanguageFile



%%%%%%%%%%%%%%%%%%%%%%%%%%%%%%%%%%%%%%%%%%%%%%%%%%%%%%%%%%%%%%%%%%%%%%

%\title{Unnecessarily Complicated Research Title}
%\subtitle{Part 1: A FAIR comic}
%\shorttitle{A FAIR comic}
%
%
% \title and \shorttitle are mandatory, \subtitle can be omitted by just leaving the string empty (please do not comment out the subtitle line)
%
%

\articletype{Research Comic}
%
%
% Select "Research Article" for manuscript submissions, "Software Descriptor" for software submissions, and "Dataset Descriptor" for dataset submissions.
%
%

\articleclass{preprint}
%
%
% To be changed by the editor:
% Enter "letter", "final" or "preprint" to determine the typesetting mode
%
%

\useScikgtex
%
%
% To be changed by author:
% Uncomment "\useScikgtex" if Scikgtex annotations are used.
% Die fuenf Pflicht-Properties stehen im Text:
%   \objective        letzter Satz der Introduction
%   \researchproblem* am Ende des Abstracts   (Sternform, siehe dort)
%   \method           im Abstract
%   \result           im Abstract
%   \conclusion*      in der Discussion, letzter Absatz (Sternform)
% Sternform = schreibt nur ins XMP, setzt keinen Text. Noetig ueberall
% dort, wo der annotierte Satz ein \cite oder \textit enthaelt.
% Dazu die bibliografischen Properties \metatitle*, \metaauthor* und
% \researchfield* direkt nach \begin{document}.
%
%

\datesubmitted{2026-05-30}
%
% State the date when you have submitted the article as a preprint.
%
%
% The information on authors and
% affiliations have to be provided in the file: authors.tex
\corresponding{\printcorrespondingauthor}
%
%
% Thanks and acknowledgements  and funding
% information have to be saved in: acknowledgements.tex
%
%

\keywords{Comic, Inggrid}
%
%
% 4-6 meaningful keywords that describe the paper



\dataavailability{Data can be found here:
\href{http://doi.org/10.5281/zenodo.19468332                                 }{http://doi.org/10.5281/zenodo.19468332}} 
\softwareavailability{-}

\addbibresource{references.bib}
%
%
% Save all your references in this Bib(La)TeX file
%
%

\pagestyle{inggrid}

\begin{document}
%
% SciKGTeX: bibliografische Metadaten des Papers. Alles Sternformen, sie
% setzen keinen Text und stehen deshalb hier gebuendelt. Erst ab
% \begin{document} verfuegbar, weil inggrid.cls das Paket per
% \AtEndPreamble laedt. Titel wie in CITATION.cff, Namen wie auf der
% Titelseite, Forschungsfeld als Label aus der ORKG-Taxonomie (R373).
\metatitle*{A FAIR Comic about Research Data Infrastructure, Part 1: From Charlemagne to ing.grid}%
\metaauthor*{Alfred Neuwald}%
\metaauthor*{Tobias Hamann}%
\metaauthor*{Évariste Demandt}%
\researchfield*{Science and Technology Studies}%
%
\maketitle
\begin{abstract}
    Comics journalism constitutes a graphic form of literary journalism. According to John S. Bak and Christopher Craig \cite{bak_comics_2026}, comics journalism and literary journalism share key features: they both read like fiction but are based on immersive research, eyewitness accounts, and verifiable facts scrutinized as closely as possible. RWTH Aachen University has recently embraced comics journalism as a means of communicating ongoing research at the Chair of AI Methodology within the Faculty of Computer Science \cite{neuwald_deus_2026}. \method{To enable close observation and immersive documentation of scientific practice, comic artist Alfred Neuwald was invited to work alongside the researchers. The resulting collaboration established the general framework for this project.} \result{Here we present the first result of a modified artist-in-residence concept implemented by the German National Research Data Infrastructure for Engineering Sciences (NFDI4ING).}%
%
% SciKGTeX: \researchproblem steht hier als Sternform. Sein Fliesstext
% enthaelt ein \cite, und der Lua-Stripper des Pakets macht daraus im XMP
% den nackten BibTeX-Key. Die Sternform setzt keinen Text, sie schreibt nur
% in die Metadaten -- der Wert ist also derselbe Satz wie oben, nur ohne
% das \cite. \method und \result bleiben direkt im Text, weil sie keine
% Zitate enthalten. \objective steht am Ende der Introduction.
\researchproblem*{RWTH Aachen University has recently embraced comics journalism as a means of communicating ongoing research at the Chair of AI Methodology within the Faculty of Computer Science.}%
\end{abstract}

\inggridlinenumbers

% Optional: normale Seite 1 ohne PDF-Hintergrund
\null

% Titel-Seite

% Disclaimer box on first FDI page
\thispagestyle{inggrid}

% !TeX root = ../main.tex


\begin{acronym}[JSONP]\itemsep0pt
\acro{aas}[AAS]{Asset Administration Shell}
\acro{ai}[AI]{Artificial Intelligence}
\acro{aims}[AIMS]{Applying interoperable metadata standards}
\acro{api}[API]{Application Programming Interface}
\acro{ardc}[ARDC]{Australian Research Data Commons}
\acro{bmbf}[BMBF]{Federal Ministry of Education and Research}
\acro{cc}[CC]{Creative Commons}
\acro{cd}[CD]{Continuous Deployment}
\acro{cern}[CERN]{Conseil Européen pour la Recherche Nucléaire}
\acro{ci}[CI]{Continuous Integration}
\acro{csv}[CSV]{Comma-Separated Values}
\acro{cpms}[CPMS]{Cyber-Physical Measurement System}
\acro{cpps}[CPPS]{Cyber-Physical Production System}
\acro{cps}[CPS]{Cyber-Physical System}
\acro{do}[DO]{Digital Object}
\acro{doi}[DOI]{Digital Object Identifier}
\acro{dfg}[DFG]{German Research Foundation}
\acro{dcc}[DCC]{Digital Curation Centre}
\acro{dlc}[DLC]{data life cycle}
\acro{dmp}[DMP]{data management plan}
\acro{dsl}[DSL]{Domain-Specific Language}
\acro{eosc}[EOSC]{European Open Science Cloud}
\acro{fact}[FACT]{Fairness, Accuracy, Confidentiality, and Transparency}
\acro{nfdi}[NFDI]{National Research Data Infrastructure}
\acro{nfdi4ing}[NFDI4Ing]{National Research Data Infrastructure for the Engineering Sciences}
\acro{fair}[FAIR]{Findable, Accessible, Interoperable, Re-usable}
\acro{fairdo}[FAIR-DO]{FAIR Digital Object}
\acro{fairse}[FAIR-SE]{FAIR Sensor Ecosystem}
\acro{fairsi}[FAIR-SI]{FAIR Sensor Infrastructure}
\acro{fairss}[FAIR-SS]{FAIR Sensor Service}
\acro{iop}[IoP]{Internet of Production}
\acro{iiot}[IIoT]{Industrial Internet of Things}
\acro{iira}[IIRA]{Industrial Internet Reference Architecture}
\acro{iot}[IoT]{Internet of Things}
\acro{lsm}[LSM]{large-scale metrology}
\acro{luh}[LUH]{Leibniz University Hannover}
\acro{mde}[MDE]{model-driven software engineering}
\acro{nist}[NIST]{National Institute of Standards and Technology}
\acro{oecd}[OECD]{Organization for Economic Cooperation and Development}
\acro{opcua}[OPC UA]{Open Platform Communications Unified Architecture}
\acro{openaire}[OpenAIRE]{Open Access Infrastructure for Research in Europe}
\acro{orcid}[ORCiD]{Open Researcher and Contributor ID}
\acro{pid}[PID]{Persistent Identifier}
\acro{prorp}[PRO\allowbreak{} Research\allowbreak{} Process\allowbreak{} 4Ing]{Project-\allowbreak{} and\allowbreak{} RDM-Oriented\allowbreak{} Research Process\allowbreak{} for\allowbreak{} Engineering Sciences}
\acro{rami}[RAMI 4.0]{Reference Architecture Model Industry 4.0}
\acro{rdaf}[RDaF]{Research Data Framework}
\acro{rdf}[RDF]{Resource Description Framework}
\acro{rdm}[RDM]{research data management}
\acro{rdmo}[RDMO]{Research Data Management Organizer}
\acro{rwth}[RWTH]{RWTH Aachen University}
\acro{smr}[SMR]{Spherically Mounted Retroreflector}
\acro{soil}[SOIL]{SensOr Interfacing Language}
\acro{ssn}[SSN]{Semantic Sensor Network Ontology}
\acro{sso}[SSO]{Single Sign-on}
\acro{tcp}[TCP]{Tool Centre Point}
\acro{tu9}[TU9]{TU9 – German Universities of Technology e.V.}
\acro{uri}[URI]{Uniform Resource Identifier}
\acro{url}[URL]{Uniform Resource Locator}
\acro{vdi}[VDI]{Association of German Engineers}
\acro{wot}[WoT]{Web of Things}
\acro{wzl}[WZL]{Laboratory for Machine Tools and Production Engineering}
\acro{wzl-iqs-en}[WZL-IQS]{Chair of Intelligence in Quality Sensing}
\end{acronym}


% \acrodef{dias}[DIAS]{dynamically interconnected assembly system}
% \acrodefplural{dias}[DIAS]{dynamically interconnected assembly systems}
% \acrodef{fjsp}[FJSP]{flexible job-shop scheduling problem} 
% \acrodef{jsp}[JSP]s{job-shop scheduling problem}

% Sprungleiste: zuerst der Artikel, dann der Comic im Fliesstext,
% danach die weiteren Sprachfassungen aus dem Anhang.
% Manuell steuerbar mit \FDIJumpTo{article,eng,deu,lat,qya}
\FDIJumpToAll

\section{Introduction}
\label{sec:introduction}
Inviting an outsider to observe and take part in scientific practice has a methodological history that can be traced back to the ethnographic approach of Bruno Latour and Steve Woolgar’s \textit{Laboratory Life} \cite{latour_laboratory_1979, latour_laboratory_1986}. \textit{Laboratory Life} is a foundational work in science and technology studies. The book includes a floor plan of the research laboratories, offices, and meeting rooms in which Latour worked alongside the researchers. It also reproduces photographs of the working environment and, in some cases, of the researchers. Latour and Woolgar thus document interactions with and among the researchers in a spatially situated manner. They report on the verbal content of these interactions and the vocabulary used. Their main focus is the researchers’ investment in the iterative production process that brings about the principal outcome of scientific work: peer-reviewed publications. Crucially, Latour and Woolgar also provide transcripts and occasionally accompany them with in-line commentary. The in-line commentary is used to describe aspects of spatially situated interaction that are not captured in the transcripts themselves, including deictic gestures---gestures that establish spatial reference within a shared environment. A key constraint of in-line commentary, however, is that it competes for space with the transcript itself: as descriptions of embodied and spatially situated action become more detailed, the transcript becomes increasingly fragmented and more difficult to read. To address this problem, multimodal forms of transcription typically move such commentary into a separate column at the page margins (see, for example, \cite{edwards_common_1987}). This layout allows verbal interaction and descriptions of embodied action to be presented simultaneously without interrupting the flow of the transcript. Moreover, this layout permits the integration of graphical representations alongside the transcript, which constitutes a comparatively recent development in conversation analytic research (see, for example, \cite{hormeyer_einsatz_2015}).

Comics represent embodied speech. They do so by depicting speakers and listeners within their environment together with spoken utterances, deictic gestures, and, in some cases, internal thoughts represented through thought balloons. As a graphic medium, comics are therefore particularly well suited to documenting multimodal interaction in situated research environments. Comics, however, usually transport readers into the realm of imagination rather than documented fact. Comics journalism---together with its European and Japanese counterparts, \textit{BD reportage} and documentary manga---differs precisely in this respect \cite{bak_comics_2026, bak_mangix_2024}.

At RWTH Aachen University, comics journalism has recently been embraced as a means of communicating ongoing research at the Chair of AI Methodology within the Faculty of Computer Science \cite{neuwald_deus_2026}. The collaboration between researchers and comic artist Alfred Neuwald established the general framework for the present project. \objective{By combining ethnographic observation with the multimodal representational capacities of comics, we explore comics journalism as a method for documenting situated research practices within the administrative office of the German National Research Data Infrastructure for Engineering Sciences (NFDI4ING).}

\section{Result}
\subsection{A FAIR Comic on Research Data Infrastructure}
% Sprungmarke des Comics im Fliesstext, passt sich \FDILanguage an
\label{comic:\FDILanguage}

\noindent
\fbox{%
  \begin{minipage}{\dimexpr\textwidth-2\fboxsep-2\fboxrule\relax}
    \textbf{Dear Reviewers, please note:} The line numbers in the following pages refer to comic panels, not to individual text lines.
  \end{minipage}%
}

\clearpage
\nolinenumbers

% Anleitung:

% \FDIText{x}{y}{breite}{text}
% x      = Abstand vom linken Seitenrand
% y      = Abstand vom oberen Seitenrand
% breite = Breite der Textbox
% text   = Inhalt der Textbox
% Beispiel: \FDIText{2cm}{3cm}{6cm}{Nordsee}

% \FDIUseText{seite}{panel}{text}
% seite = Nummer der PDF-Seite
% panel = Nummer des Panels auf dieser Seite
% text  = Nummer des Textes in diesem Panel
% Beispiel: \FDIUseText{1}{1}{1}

% \comic{sprache}
% gibt alle Comicseiten der gewaehlten Sprache direkt hintereinander aus
% sprache = ISO-639-3-Code des Ordners in language_files,
% z.B. deu, eng, lat, qya
% Beispiel: \comic{deu}
% Optional laesst sich die Seitenzahl begrenzen: \comic[3]{deu}
% Mehrere Sprachen nacheinander sind moeglich, z.B. \comic{deu}\comic{eng}

% Einzelne Seiten koennen weiterhin mit \FDIInputPage{seite} in der
% Dokumentsprache \FDILanguage ausgegeben werden.

\comic{\FDILanguage}

\clearpage
\pagestyle{inggrid}
\inggridlinenumbers

\section{Materials and Method}
\subsection{Project Development and Collaboration}
Évariste Demandt first contacted Alfred Neuwald in July 2025 and met with him shortly thereafter. During this initial meeting, he learned about Alfred's ongoing collaboration with the Chair of AI Methodology. As this collaboration was expected to continue for several more months, Alfred's availability for additional commitments was limited. Nevertheless, both parties agreed to remain in contact and continue exchanging ideas. 
Évariste suggested that Alfred might participate as a guest observer in the 2025 Conference on Research Data Infrastructure (CoRDI 2025), and Tobias Hamann subsequently extended a formal invitation. 
CoRDI 2025 took place at RWTH Aachen University from 26 to 28 August 2025 \cite{nfdi_ev_cordi_nodate}. 
Alfred participated in the
%first day of the 
conference, including the evening event at the \ac{wzl}.
He used the opportunity to familiarize himself with the field of research data infrastructure and with the people involved.

Following the conference, Évariste convened a meeting between Alfred and representatives of NFDI4ING, IT Center and \ac{wzl}, as well as a representative of the Chair of Society and Technology at RWTH Aachen University. It quickly became apparent that we lacked the resources to replicate Alfred's ongoing collaboration with the Chair of AI Methodology, a project that enabled Alfred to work full-time with the researchers. As a result, we began exploring alternative formats for collaboration. Recognizing Alfred's ability to ask questions, Évariste concluded that a format centered on regular conversations would best accommodate the budget constraints. A weekly \textit{jour fixe} was therefore scheduled, during which Alfred visited NFDI4ING's administrative office at the IT Center.

Alfred began work on the FAIR comic in February 2026. He photographed Évariste and, later, Tobias as reference material for the comic, which he created outside RWTH Aachen University at his atelier. The collaboration followed a meeting-only format: no emails, phone calls, or messages were exchanged between meetings in order to keep the workload manageable. This meant that Alfred always came in well prepared with questions and, later, drafts. Meetings were open-ended and typically exceeded one hour. Alfred's questions sometimes revealed gaps in shared understanding of \ac{rdm} or exposed issues that could not be resolved immediately. In such instances, the questions were adjourned and additional expertise was sought from specialists. This occurred twice during the project: once with support from NFDI4ING's Task Area Federated Interactive Data Space and once with support from the Chief Technical Officer of the IT Center.

During the weekly meetings, learning was reciprocal. Alfred shared his extensive knowledge of comic culture, referring to documentary works such as Art Spiegelman's \textit{Maus} \cite{spiegelman_maus_1986}, Marjane Satrapi’s \textit{Persepolis} \cite{satrapi_persepolis_2003}, and Guy Delisle's \textit{Pyongyang} \cite{delisle_pyongyang_2003}. Delisle's work, in particular, received special attention as an outstanding example of observer commentary. Alfred admired Delisle's ability to document a social world and occasionally drew parallels between his own encounters with researchers and Delisle's experiences working in North Korea.

\subsection{Documentary Basis and Narrative Framing}
This is a work of literary journalism, which means that it is grounded in Alfred's research. The FAIR comic conveys his first-hand observations, conversations, and photographic documentation gathered during the weekly meetings at the administrative office of NFDI4ING at the IT Center of RWTH Aachen University. Among the materials that informed Alfred's work was a leaflet published by the Landesinitiative NFDI der Digitalen Hochschule NRW and the AG FDM Awareness \cite{landesinitiative_nfdi_der_digitalen_hochschule_nrw_nachnutzbare_2019}. Copies were available in the administrative office, and the leaflet became a recurring point of reference in conversations between Évariste and Alfred and, later, between Évariste, Tobias, and Alfred. The influence on the FAIR comic can be seen both in the range of subjects covered and in the canned-food metaphor, which originates on page 12 of the leaflet. Another material used by Afred was a slide desk from NFDI4ING's Task Area Federated Interactive Data Space with a schematic representation of the different layers that constitute a FAIR Digital Object.

A point of contention was whether fictional elements should be included at all. At the outset of the project, Évariste was unfamiliar with the distinction between fictional comics and comics journalism and proposed narrative devices drawn from fictional comics. Alfred generally resisted such proposals. The need to reconcile them with the request to begin the story in 843 AD ultimately led to the inclusion of a few fictional elements on the first three pages, which are examined below. Before turning to these elements, however, we first explain Évariste's request to begin the story in 843 AD.

At the time of writing, Aachen has presented the International Charlemagne Prize for seventy-six years. First awarded on 14 March 1950, the prize honours individuals who have made outstanding contributions to peace, unity, and cooperation in Europe \cite{stiftung_internationaler_karlspreis_zu_aachen_start_nodate}. Its creation coincided with a formative decade of European integration. The decade witnessed the establishment of the European Coal and Steel Community (ECSC) through the Treaty of Paris (1951; in force from 1952) and, later, of the European Economic Community (EEC) and the European Atomic Energy Community (EURATOM) through the Treaties of Rome (1957; in force from 1958). Less frequently noted is that the combined territory of the six states that founded the ECSC, the EEC, and EURATOM overlaps substantially with the core regions of the Carolingian Empire. This overlap is noteworthy because Aachen is closely associated with Charlemagne, who established a palace complex (Pfalz) there and made the city a center of Carolingian rule. Aachen therefore occupies a distinctive place in both the symbolic landscape of post-war European integration and the historical geography of the Carolingian world.

The Carolingian connection is also visible at RWTH Aachen University. It is reflected, among other things, in the acronym C.A.R.L., the name of a representative university building that served as the venue for CoRDI 2025 \cite{nfdi_ev_cordi_nodate}. Moreover, the local host of CoRDI 2025, Robert Schmitt of NFDI4ING, referred explicitly to the Carolingians in his opening remarks of the CoRDI. Against this backdrop of Carolingian references, Évariste was reminded of a familiar narrative convention in comics: the use of an introductory map, as found, for example, on the front endpapers of \textit{Astérix} albums \cite{goscinny_asterix_1961}. This convention inspired the idea of opening the FAIR comic with a map of the Carolingian Empire.

The first three pages of the FAIR comic are dedicated to establishing context and continuity (cf. chapter five in \cite{edwards_common_1987}) within the broader geographical and cultural landscape of NFDI4ING. They are also intended to entertain readers and draw them toward the more technical aspects of research data infrastructure. To this end, a small number of fictional elements were introduced, including reconstructions of what Charlemagne and his descendants might have said and done up to the partition of the empire under the Treaty of Verdun in 843 AD. In designing and drawing the first three pages, Alfred used expertise he had acquired in the course of an earlier project \cite{neuwald_karl_2014}. The map of the Carolingian Empire, specifically, is based on a historical map distributed through Alamy \cite{noauthor_europe843_nodate}. On a more technical note, most of the details of the historical map were omitted in order to emphasize selected features, while a few others were added. Among these is the empire's eastern frontier, which is marked on the map by the added detail \textit{limes sorabicus}. The territories beyond this frontier were subject to varying degrees of Carolingian rule and correspond broadly to the toponomastic boundary between regions characterized by Germanic and Slavic place-name strata \cite{jelic_ethnonym_2019}. Another addition is the city of Darmstadt, home to Technische Universität Darmstadt, applicant institution of NFDI4ING during its second funding period \cite{noauthor_dfg_nodate}.

As soon as the FAIR comic shifts from the Carolingian past to present-day Aachen and C.A.R.L., the staged setting of the meeting between Alfred and Évariste, the colour palette switches to indicate a different register: conversation. The ensuing dialogue depicts discussions that took place over the course of almost five months of weekly meetings involving Alfred and either Évariste alone or Évariste and Tobias together. The FAIR comic organizes these discussions thematically and condenses them into nine pages.

\subsection{Manuscript Preparation and Translation}
The process of outlining, designing, and refining the FAIR comic was conducted in German. Accordingly, the first draft was written in German. As we approached publication of the FAIR comic in \textit{ing.grid}, however, it had to be translated into English. At the same time, we considered making the FAIR comic available in different languages important for its accessibility. One way of doing so would have been for Alfred to redraw all text in the FAIR comic for each additional language. This would have resulted in lengthy review cycles and placed a considerable additional workload on Alfred.

Tobias instead developed a more automated approach to producing multilingual versions. Alfred created a blank PDF version of the FAIR comic containing the interlocutors, maps, and all other graphical elements, but no text. This version served as the basis for a {\LaTeX} document in which each blank comic page was used as the background of a full-page {\TikZ} image \cite{Tantau_2013_book}. {\TikZ} is a recursive acronym for \textit{{\TikZ} ist kein Zeichenprogramm}, German for \enquote{{\TikZ} is not a drawing programme}. Alfred's original German text was then superimposed on the pages using {\TikZ} nodes (see the \hyperref[comic:deu]{German version in the appendix}). These nodes can be positioned freely and styled individually, allowing precise control over fonts, font sizes, and special text arrangements, including curved or rotated text. Manual line breaks can likewise be specified in the language file, allowing the text layout to be adjusted to achieve the visual balance intended by the artist. One language-dependent element, however, could not be handled as a textual overlay. Alfred also had to remove the drawing of three-dimensional toy bricks forming the letters \enquote{FDI}, standing for the German \textit{Forschungsdateninfrastruktur}, from the first page. Here, a language-dependent version of the image was required. To address this, different language versions of the image were generated using artificial intelligence (AI), namely OpenAI's ChatGPT 5.5. The \ac{ai} was instructed to preserve the style of the original image while replacing the hand-sketched \enquote{FDI} with \enquote{RDI}.

Similarly, the language file can be translated automatically using \ac{ai}. To test this concept, we had \ac{ai} translate the language file into Latin and J. R. R. Tolkien's Elvish, thereby generating corresponding versions of the FAIR comic. Both versions have a dextrograde (left-to-right) writing direction. The Elvish version, rendered in the Tengwar script, further demonstrates that non-Latin scripts can be accommodated. Sinistrograde (right-to-left) and boustrophedon (alternating writing direction from line to line) writing have not yet been implemented or tested.

The source code is available on \href{https://github.com/tobias-hamann/research-comic-vol1}{GitHub}\footnote{\url{https://github.com/tobias-hamann/research-comic-vol1}} \cite{Hamann_A_FAIR_Comic_2026}. 
Every language \verb|CODE| is kept in its own folder \path{language_files/CODE}, which holds the text file \path{CODE.tex}, the corresponding folder \path{pages_CODE} and all language dependent images. 
Here, \verb|CODE| is the lowercase three-letter identifier from ISO 639 Set 3 (formerly ISO 639-3) \cite{iso_639_2023}, as used by the Zenodo language vocabulary.
The Draconic version uses the project-local placeholder \verb|vss|, derived from the Draconic name \emph{Vs'shtak}; this identifier is neither registered in ISO 639 nor in Zenodo.
If additional fonts are needed, they are contained in the subfolder \path{fonts_CODE}.
As a wrapper, the file \path{comic_files/pages_functions.tex} is processing these folders.
Every \path{language_files/CODE/CODE.tex} contains all text lines ordered by their appearance on the pages in the corresponding language, which allows for a translation of texts independently from their appearance in the document.
Also, manual line breaks can be added here as needed to account for the styling and weighting of texts within speach bubbles.
The \path{pages_CODE} files handle the positioning and scaling of texts for each language individually, which accounts for different lengths of the same contents in different languages.
The language is chosen in the \verb|\comic{CODE}| command.
It outputs all comic pages in the selected language one after another, e.g. \verb|\comic{eng}| generates the English version of the comic.
Since each call is self-contained, several language versions can also be printed one after another within the same document.
The primary comic language and the language-specific title, subtitle, and short title are selected in \verb|main.tex| with \verb|\newcommand{\FDILanguage}{eng}|. The article body and the remaining metadata are maintained separately.
As of writing, the original \hyperref[comic:deu]{German} version and this translated \hyperref[comic:eng]{English} version exist along the Latin, Elvish, and Draconic versions.

\section{Discussion}
\subsection{Dialogue as a Documentary and Didactic Form}
The conversational format of the FAIR comic has a philosophical pedigree reaching back to the dialogues of Plato (\textit{circa} 424/3--348/7 BC \cite{debra_people_2002}). Parallels include the use of historical interlocutors \cite{debra_people_2002} and a characteristic division of labour: Socrates asks questions, while his interlocutors formulate and defend answers. Like Socrates, Alfred frequently advances the discussion through questions. One of the most arresting things about the Platonic dialogues is that most of them have survived for more than two millennia and continue to be widely read and discussed today. We contend that such an unusually long period of transmission should be taken as an indication of the longevity not only of their content but also of their form. References to Plato continue to appear in contemporary culture and education. Examples include Jay-Z, Kanye West, Frank Ocean, and The-Dream's song "No Church in the Wild", which invokes the \textit{Euthyphron} dilemma \cite{jay-z_no_2011, kanye_west_jay-z_2012}; Helme Heine's children’s book \textit{Tante Nudel, Onkel Ruhe und Herr Schlau}, which adapts the \textit{Politeia} \cite{heine_1979}; Höfner and Süßbier's didactic comic \textit{Das verrückte Mathe-Comic-Buch}, which includes a re-telling of the \textit{Menon} \cite{hofner_verruckte_2012}; as well as James Orr's lecture series \textit{Plato} \cite{peterson_academy_dr_2026}. This continued presence suggests not only the enduring relevance of the subject matter, but also the effectiveness and remarkable adaptability of the dialogical format through which that subject matter is explored. As Orr points out in the opening lecture of his series, "once you know all of the sides, you understand your own" \parencite[at 55:09]{peterson_academy_dr_2026}. The FAIR comic similarly relies on dialogue as a means of exploring and communicating complex, albeit technological, matters. We hope that readers interested in research data infrastructure will be profitably entertained by the result.

We have introduced the FAIR comic by reference to the documentary practices of ethnographic research \cite{latour_laboratory_1979, latour_laboratory_1986}, and discourse and conversation analysis \cite{edwards_common_1987, hormeyer_einsatz_2015}. More specifically, we have situated comic culture within a continuum of practices concerned with documenting deixis and embodied speech (see \ref{sec:introduction}). We have argued that, as a graphic medium, comics are particularly well suited to documenting multimodal interaction in situated research environments. In this sense, the FAIR comic may itself be understood as a documentary form of speech. If we have now reached back to the dialogues of Plato, it is still in an attempt to reflect on the structure of the FAIR comic and, more generally, on what helps readers understand the spoken word \cite{ong_orality_1982}.


\subsection{The FAIR Comic between Research and Science Comics}
We conclude furthermore our discussion by comparing the FAIR comic to related forms of comic culture, such as comics about researchers in research institutions (see, for example, Jorge Cham's \textit{PHD Comics} or Eppendorf’s \textit{Lab Life}; \cite{cham_phd_1997, noauthor_lab_nodate}), as well as science communication comics that explain fields of research (see, for example, Larry Gonick's \textit{Cartoon Guide} series or Matteo Farinella's \textit{Neurocomic}; \cite{gonick_cartoon_1983-1, gonick_cartoon_1983, farinelli_neurocomic_2013}). The FAIR comic sits comfortably between these two broader lineages and follows a journalistic ethos. What we seek to achieve is not adequately captured by saying that the FAIR comic is about research practice, nor even that it is researched, but rather that it is explicit about the process of learning---or coming to understand. To put our claim in direct speech: "This is what we discussed while trying to establish a mutual understanding of research data infrastructure."%
%
% SciKGTeX: \conclusion als Sternform, gleiche Begruendung wie im Abstract.
% Hier kommen zu den \cite noch die verschachtelten \textit{..} dazu, an
% denen der Stripper verwaiste Klammern hinterlaesst. Der Wert unten ist
% der Absatz oben ohne Zitate, ohne Kursivierung, mit echtem Geviertstrich.
\conclusion*{We conclude furthermore our discussion by comparing the FAIR comic to related forms of comic culture, such as comics about researchers in research institutions (see, for example, Jorge Cham's PHD Comics or Eppendorf’s Lab Life), as well as science communication comics that explain fields of research (see, for example, Larry Gonick's Cartoon Guide series or Matteo Farinella's Neurocomic). The FAIR comic sits comfortably between these two broader lineages and follows a journalistic ethos. What we seek to achieve is not adequately captured by saying that the FAIR comic is about research practice, nor even that it is researched, but rather that it is explicit about the process of learning—or coming to understand. To put our claim in direct speech: "This is what we discussed while trying to establish a mutual understanding of research data infrastructure."}%


%
%
% Part of SciKGTex
%
% Die fuenf Pflicht-Properties sind beim zugehoerigen Fliesstext gesetzt,
% siehe die Uebersicht oben bei \useScikgtex. Weitere Properties gehen mit
% \addmetaproperty{name} und \contribution*{name}{wert}, wobei die
% Sternform den Wert nur in die Metadaten schreibt, ohne ihn zu setzen.
%
%

%
%
% Appendix: alle weiteren Sprachfassungen des Comics
%
% \FDIAppendixComics gibt jede Sprache aus \FDILanguageList aus, ausser
% der Sprache, die schon im Fliesstext steht. Jede Fassung bekommt eine
% Ueberschrift und die Sprungmarke comic:sprache.
% Reihenfolge und Auswahl stehen in \FDILanguageList in
% comic_files/pages_functions.tex, aktuell: deu, eng, vss.
% Einzeln geht auch: \FDIComicSection{deu} oder \FDIComicSection[3-5]{deu}
%
\appendix
\FDIAppendixComics

%
%
% Acknowledgements, contributions & roles, and references
% will be set automatically at the end of the document
%
%
\end{document}
